\documentclass[11pt,a4paper]{article}
\usepackage[margin=25mm]{geometry}
\usepackage{amsmath,amssymb,amsthm,booktabs,hyperref}
\setlength{\emergencystretch}{2em}
\newtheorem{theorem}{Theorem}
\newtheorem{lemma}[theorem]{Lemma}
\newcommand{\Aut}{\operatorname{Aut}}
\title{Bounds for generating pairs from maximal subgroups}
\author{Human Compression project: a component of the L046 audit}
\date{}
\begin{document}\maketitle
\begin{abstract}
A union bound over maximal subgroups gives a lower bound for the number
of generating-pair orbits, without the full subgroup lattice. Applied to
the O'Nan group with the published maximal-subgroup classification and
ordinary character tables, it gives at least 798370 orbits in the
convention of Connor and Leemans. This exceeds the total 796772 in
their Theorem 1.1(2). One table entry can be determined exactly:
the type $(20,31)$ has 9118 orbits, rather than the printed 7246.
The same calculation proves existence for all 129 types asserted to
exist in that paper. We do not determine the exact total or the full
chirality classification.
\end{abstract}

\section{The counting convention and the bound}
For a finite group $G$, let
\[
 \begin{gathered}
 X_G(p,q)=\{(R,S):|R|=p,\ |S|=q,\ |RS|=2\},\\
 E_G(p,q)=\{(R,S)\in X_G(p,q):\langle R,S\rangle=G\}.
 \end{gathered}
\]
Write their cardinalities as $x_G(p,q)$ and $e_G(p,q)$, and write
$a=|\Aut(G)|$. Automorphisms act freely on $E_G(p,q)$: fixing both
generators fixes every element of $G$. Consequently the number of
orbits is $n_G(p,q)=e_G(p,q)/a$.

Pairs are ordered. Taking $p\le q$ selects one orientation of each
type and makes no further duality identification when $p=q$.
This is precisely the source's convention in its product-one triples
and formula (2.4). Although the source also writes braces for pairs,
the result below concerns the count defined by that formula and
reported in Table 1 and Theorem 1.1(2).

\begin{theorem}\label{union}
Let $M_1,\ldots,M_s$ represent all conjugacy classes of maximal proper
subgroups of $G$. Put
\[
 B_G(p,q)=\sum_{j=1}^s [G:M_j]x_{M_j}(p,q).
\]
Then
\[
 \max\left(0,\left\lceil\frac{x_G(p,q)-B_G(p,q)}a\right\rceil\right)
 \ \le\ n_G(p,q)\ \le\ \left\lfloor\frac{x_G(p,q)}a\right\rfloor. \tag{1}
\]
\end{theorem}
\begin{proof}
Every pair that does not generate $G$ is contained in a maximal proper
subgroup, because $G$ is finite. A conjugate of $M_j$ contains precisely
$x_{M_j}(p,q)$ pairs of the prescribed type. The number of its conjugates
is $[G:N_G(M_j)]$, which is at most $[G:M_j]$. Taking the union of these
sets, and allowing repeated counting in intersections, gives
\[
 |X_G(p,q)\setminus E_G(p,q)|\le B_G(p,q).
\]
Subtract from $x_G(p,q)$ and divide by $a$. The number of generating
orbits is a nonnegative integer, giving the lower bound. The inclusion
$E_G(p,q)\subseteq X_G(p,q)$ gives the upper bound, again using
integrality. No assumption that maximal subgroups intersect trivially
is made. Repeated isomorphism types in the list must retain their
conjugacy-class multiplicities.
\end{proof}

The character calculation uses the Frobenius identity already acquired
in the separately reviewed manuscript \emph{Generating pairs from
characters and subgroup incidence}. If $C_i$ are the conjugacy classes
of $H$, define $v_{\chi,t}=\sum_{|c_i|=t}|C_i|\chi(c_i)$. Then
\[
 x_H(p,q)=\frac1{|H|}\sum_{\chi\in\operatorname{Irr}(H)}
             \frac{v_{\chi,2}v_{\chi,p}v_{\chi,q}}{\chi(1)}. \tag{2}
\]
The acquired proof explains the bijection with product-one triples.
All arithmetic in (2) is exact. In the retained tables every sum
$v_{\chi,t}$ is an integer, verified before export. A separate Python
program then evaluates (2) with rational numbers from these integer
vectors. Neither source generating-pair counts nor expected totals
are inputs to (2) or (1).

\section{The retained subgroup input}
Let $G$ be the O'Nan sporadic simple group. We retain its established
order $460815505920$, the identification $\Aut(G)=G:2$, and the complete
classification of its thirteen conjugacy classes of maximal subgroups.
The following table records their isomorphism types, orders and
multiplicities. The names in the last column identify the actual
CTblLib tables used; distinct names can describe isomorphic groups.
\begin{center}\small
\begin{tabular}{lrrl}\toprule
type&order&classes&character tables\\\midrule
$L_3(7):2$&3753792&2&\texttt{L3(7).2}, \texttt{ONM2}\\
$J_1$&175560&1&\texttt{J1}\\
$4_2.L_3(4):2_1$&161280&1&\texttt{4\_2.L3(4).2\_1}\\
$(3^2:4\mathbin\times A_6).2$&25920&1&\texttt{ONM5}\\
$3^4:2^{1+4}.D_{10}$&25920&1&\texttt{3\^{}4:2\^{}(1+4)D10}\\
$L_2(31)$&14880&2&\texttt{L2(31)}, \texttt{ONM8}\\
$4^3.L_3(2)$&10752&1&\texttt{4\^{}3.L3(2)}\\
$M_{11}$&7920&2&\texttt{M11}, \texttt{ONM11}\\
$A_7$&2520&2&\texttt{A7}, twice\\\bottomrule
\end{tabular}\end{center}
The product notation for the order-25920 subgroup follows the online
Atlas display; the exact named character table, not this abbreviation,
is the finite input used in the computation.

These are external classification premises. We acquire the complete
list from the Atlas of Group Representations, cross-checked against
CTblLib's stored maximal-class list. Wilson's survey, Section 3.3,
attributes independent determinations to Yoshiara, Ivanov--Tsaranov--
Shpectorov and Wilson. We do not reproduce those classification proofs.
The named character tables are likewise retained Atlas/CTblLib data;
checking their orders or character identities alone would not prove
their group identifications. The missing full subgroup-incidence matrix
of $G$ is not required for the weaker union bound (1).

\section{The two numerical conclusions}
Use exactly the source's set of rotation orders
\[
 P=\{3,4,5,6,7,8,10,11,12,14,15,16,19,20,28,31\}.
\]
For each of its 136 pairs $p\le q$, let $l(p,q)$ be the lower bound
in (1). The complete ledger records $x_G$, each of the thirteen terms
of $B_G$, and both integer bounds. The sums of $l(p,q)$ with a fixed
first coordinate are as follows.
\begin{center}\begin{tabular}{r|r@{\qquad}r|r}\toprule
$p$&row lower sum&$p$&row lower sum\\\midrule
3&274&12&29404\\
4&5124&14&36418\\
5&7712&15&43644\\
6&18456&16&228833\\
7&29128&19&88545\\
8&84204&20&33646\\
10&63643&28&13903\\
11&109437&31&5999\\\bottomrule
\end{tabular}\end{center}

\begin{theorem}\label{total}
In the source's counting convention,
\[
 \sum_{p\le q\in P} n_G(p,q)\ \ge\ 798370\ >\ 796772.
\]
Thus Theorem 1.1(2) of the fixed published source is false.
\end{theorem}
\begin{proof}
The types are invariant under automorphisms, so their orbit sets are
disjoint. Sum (1) over the 136 types. The sixteen displayed integer
row sums add to 798370. The source asserts the smaller total 796772.
The discrepancy is 1598 even with this lower bound; we make no claim
that the lower bound is the exact total.
\end{proof}

\begin{theorem}\label{entry}
The exact value for type $(20,31)$ is $n_G(20,31)=9118$.
The corresponding entry of the source's Table 1 is 7246.
\end{theorem}
\begin{proof}
Among the listed maximal subgroups only the two classes of $L_2(31)$
can contain an element of order 31. Their element-order set is
$\{1,2,3,4,5,8,15,16,31\}$, which excludes 20. This set is read from
the retained character tables. It also follows from the usual split,
nonsplit and unipotent alternatives in $\operatorname{PSL}_2(31)$,
whose nontrivial orders divide 15 or 16, or equal 31.
Thus no maximal subgroup contains elements of both prescribed orders,
and every pair in $X_G(20,31)$ generates $G$. Equation (2) gives
\[
 x_G(20,31)=8403431565957120,
 \qquad |\Aut(G)|=921631011840.
\]
Their quotient is 9118. The printed entry 7246 was checked directly
in the native rendering of source page 11, in addition to extraction.
\end{proof}

Exactly 129 lower bounds in the ledger are positive. Their types are
precisely the source's asserted positive types, so their existence is
recovered simultaneously. The seven types with zero lower bound are
\[
 (3,3),(3,4),(3,5),(3,6),(3,7),(3,12),(4,4).
\]
A zero lower bound is not a nonexistence proof. This component does
not infer the exclusion of all seven. Twenty of the 136 intervals have
equal endpoints, which supplies exact counts for those types; the
other intervals are kept as intervals. Apart from $(20,31)$, every
printed entry lies within its calculated interval. That consistency
is not a verification of the other 116 unresolved exact counts.

\section{Verification, attribution and limits}
The executable certificate freshly extracts all character vectors from
the named tables and recomputes every rational sum and integer bound.
Repeated maximal classes, including the two copies of $A_7$, are
preserved. The source table is used only afterward for comparison.
Separate complete controls in $S_3,D_8,A_4,S_4,A_5$ enumerate all
subgroups and all ordered pairs of each type and check the union
inequality and free-orbit divisibility literally.

The proof uses the standard union bound, ordinary character theory
and established subgroup classifications. The contribution is their
specific application and checked numerical consequence for the fixed
source. We claim neither a new maximal-subgroup classification nor
priority for any general method. This is a correction to a published
assertion, not the settlement of a conjecture. The exact full total,
the remaining Table 1 entries and the chirality and inversion spectra
remain separate obligations. No Lean verification or numerical
compression factor is supplied by this certificate.

\begin{thebibliography}{9}
\bibitem{CL} T. Connor and D. Leemans, \emph{Algorithmic enumeration of
regular maps}, Ars Mathematica Contemporanea 10 (2016), 211--222.
\bibitem{Wilson} R. A. Wilson, \emph{The maximal subgroups of the O'Nan
group}, Journal of Algebra 97 (1985), 467--473.
\bibitem{survey} R. A. Wilson, \emph{Maximal subgroups of sporadic groups},
survey dated 8 April 2016, Section 3.3.
\bibitem{Atlas} R. A. Wilson, S. J. Nickerson, J. N. Bray et al.,
\emph{Atlas of Group Representations}, O'Nan group page,
\url{https://brauer.maths.qmul.ac.uk/Atlas/v3/spor/ON/},
accessed 13 October 2026.
\end{thebibliography}
\end{document}
